\section{从完整连接锥到全局 Cauchy 发展}\label{sec:global}
\subsection{外部定理的使用范围}
本节使用\cite{KU}的经典局部适定性和几何拼接结果。为了明确证明依赖，列出实际调用的接口：
\begin{enumerate}
\item Definition 2.3 及 Propositions 2.3--2.4：完整球面数据与约束输运；
\item Proposition 3.1：给定两个模型时空矩形和其间的完整特征连接，得到包含它们的局部时空拼接；
\item Lemmas 4.1--4.2：Minkowski 和 ERN 球面数据的规范轨道识别；
\item §5.2 及其脚注 7：把连接锥放入正规中心、单端渐近平直 Cauchy 布局；中心延拓损失一个导数。Lemma 5.1 给出以 Cauchy 稳定性实现同一步骤的方法；
\item Appendix B, Proposition B.2：在相应球对称全局框架中，由膨胀符号得到对一般闭曲面的陷获或反陷获排除结论。
\end{enumerate}
不直接把带大质量条件的 Corollary 1 应用于小耦合。本文替换的是其中 Theorem 2B 产生连接锥的输入，后续局部存在与全局几何布置按其证明实施。下文说明新输入满足那些步骤所用的条件。

\subsection{两端模型及平直中心}
给定目标 $k\ge2$，取 $K=k+8$ 并应用定理\ref{thm:cone}。左端记 $r=\rho>0$、$r_t=\beta>0$，$r_U=-1/(4\beta)$。选择一个真实 Minkowski 双零矩形，把这一球面置于所需角点；右端选择 ERN 外部及跨视界局部矩形，把 $r=Q=1$ 的终端球面置于对应角点。两端的完整数据分别在其模型规范轨道内，所以 Proposition 3.1 给出两模型之间的局部 $C^K$ 拼接。

左端自由标量在开区间恒零，且初始全部数据为 Minkowski。通过局部唯一性，这一段附近是实际平直开区域。在同一仿射规范中，锥向过去的平直延长为
\begin{equation}\label{eq:centerextension}
r(t)=\rho+\beta(t-t_-),\qquad
t_c=t_- -\rho/\beta.
\end{equation}
在 $t_c$ 达到正则中心。每个固定脉冲成员的 $\rho,\beta$ 都有限且严格正，所以可选一个包含该顶点的平直区域。这里的几何大小依赖成员即可。§5.2 示意构造中的固定起始半径可换成 $\rho$：平直锥的坐标平移与正的零坐标缩放改变坐标位置和斜率，不改变耦合 $\mu$。

\subsection{中心延拓与正则性余量}
球对称局部存在区域可按\cite[§5.2]{KU}引用的正规中心延拓步骤延至中心。其脚注明确指出，相对于特征数据，中心处波方程损失一个导数。本稿采用 $K=k+8$，足以覆盖这一损失及从时空取类空初值时的一阶损失。

也可用其 Lemma 5.1 检查此处并不需要脉冲族的一致中心邻域：从已有局部区域选择一块类空初值，内球精确平直，稍外的环由真实局部解诱导。数据满足约束。在平直球边界，各差分的低阶导数消失；对宽度 $\epsilon$ 的外环截断，若初值至少 $C^5$，Leibniz 公式与 Taylor 余项给截断差分的 $H^4$ 范数为 $O(\epsilon)$。固定当前脉冲后取小 $\epsilon$，Cauchy 稳定性给覆盖整个平直锥顶点的一个正宽度发展域。辅助截断远处允许不满足约束，只取原约束数据的依赖域；约束传播使这一依赖域内为原 EMCSF 解。高阶正则性保持给所需有限 $C^k$ 解。$K\ge10$ 保证上述有限阶要求。

\subsection{单端初始面、外部完备性与事件视界}
在上述拼接几何中选一条球对称类空曲线 $\Sigma$：从顶点之前的平直中心出发，中间紧贴连接区域的外侧，远端接入 ERN 的标准渐近平直空间端，例如在该端取静态坐标中的常时间曲线。旋转后 $\Sigma\cong\R^3$，具有一个渐近平直端。其数据从真实时空诱导，所以满足约束。

连接锥整段 $r_U<0$，紧段上有逐成员的严格余量。因而 $\Sigma$ 在连接区可选得足够近，使入射膨胀仍负；平直中心和 ERN 外段同样有这一符号。Raychaudhuri 方程在最大未来发展中传播负入射膨胀。

右侧精确 ERN 外部由相容的有限双零矩形穷尽，局部唯一性使它们在重叠处一致。靠近未来视界内侧的存在邻域允许随先进时间缩小，所需的是覆盖每个有限视界点的开邻域，而非一张具有统一宽度的无限矩形。精确 ERN 外部提供完整的 $\mathcal I^+$。

采用§5.2的因果布局，在球对称商空间中把连接锥定为 $U=0$，外侧为 $U<0$。外侧的外向零线可以进入既定 ERN 外部并抵达该端 $\mathcal I^+$；内侧未来因果曲线的两零坐标单调，不能穿回 $U<0$ 逃到这个外部。把平直中心延长段接入后，这一分界为
\[
\mathcal H^+=\partial J^-(\mathcal I^+).
\]
视界内侧有局部开放邻域，故黑洞区非空。$\Sigma$ 选择在视界之前，且全部位于可逃逸一侧。

最后取 $\Sigma$ 诱导数据的唯一最大未来整体双曲发展。由依赖域与唯一性，它包含拼接时空中 $\Sigma$ 未来的依赖域，特别包含上述外部及视界。晚时外通信域连同视界因此精确为 ERN。无反陷获对称球通过 Raychaudhuri 传播；在此正规单端球对称框架内，Proposition B.2 进而排除任意反陷获闭曲面，并排除外通信域中的陷获闭曲面。因此初始 $\Sigma$ 无陷获或反陷获闭曲面。

这些全局步骤所需的是正半径、严格负入射膨胀、完整有限阶匹配以及两端精确模型区域；没有再次使用原大质量构造的定量不等式。第\ref{sec:positive}--\ref{sec:lift}节恰好提供了这些条件。

\subsection{尺度恢复与主定理的结论}
当前解的归一化质量为一、耦合 $e=\mu$。固定坐标下作
\begin{equation}\label{eq:scaling}
g\mapsto M^2g,\quad A\mapsto MA,\quad F\mapsto MF,
\quad\phi\mapsto\phi,\quad e\mapsto e/M.
\end{equation}
约化方程或四维场方程直接验证这一缩放保持 EMCSF，质量、电荷和面积半径都乘 $M$。于是物理耦合为 $|\mathfrak e|=\mu/M$，无量纲乘积保持 $\mu$。复共轭和电荷取向变换处理另一符号。

对每个固定 $\mu>1/2$、有限 $k\ge2$，上述构造给定义\ref{def:class}中的一个解，所以 $(1/2,\infty)\subset\mathscr A_k$。定理\ref{thm:necessary}给反向包含，故
\[
\boxed{\mathscr A_k=(1/2,\infty).}
\]
同时 $1/2$ 不可达，完成定理\ref{thm:main}的证明。

\subsection{证明边界与可复核范围}
本文的主结论依赖\cite{KU}已建立的 PDE 局部适定性、中心延拓和全局因果框架；这些外部定理并非在此从头重证。内部证明的关键可核查节点是\eqref{eq:square}、\eqref{eq:compactrec}--\eqref{eq:KODE}、\eqref{eq:diagonal}、\eqref{eq:infinite-energy}--\eqref{eq:Hcomparison}以及\eqref{eq:Xtransport}--\eqref{eq:jacobian}。

已有审计包的两套符号脚本分别报告 26 项和 18 项代数检查通过，并在普通与优化执行模式下取得相同结果；其中独立 Volterra 核递推检查到第十阶。这些记录只支持相应代数，不认证鞍点分析、无限区间误差或全局 PDE 证明。全阶公式的证明在命题\ref{prop:diagonal}，并不由十阶检查外推。

本文不提供单个 $C^\infty$ 成员。实际上，固定有界非零剖面若满足所有 $M_s(z)=0$，则有限复测度 $z(-\log t)\dd t$ 消去所有多项式；由多项式在 $C([0,1])$ 的稠密性，该测度只能为零。因此不能把当前有限阶修复对同一剖面直接令 $K\to\infty$。这只是当前方法的量词边界，不是非线性光滑拼接不可能性定理。
